\documentclass[a4paper,11pt]{article}
\usepackage[utf8]{inputenc}
\usepackage[spanish]{babel}
\usepackage{amsmath}
\usepackage{amssymb}
\usepackage{geometry}
\usepackage{fancyhdr}
\usepackage{siunitx}

\geometry{margin=2cm}
\pagestyle{fancy}
\fancyhf{}
\rhead{Examen 2022 OPTIQUE - Soluciones}
\lhead{AIST}
\rfoot{\thepage}

\title{\textbf{Examen 2022 OPTIQUE - Soluciones\\Submarine Cable Systems}}
\author{}
\date{}

\begin{document}
\maketitle

\section*{Preguntas Teóricas}

\subsection*{1) Fórmula Shannon-Hartley}
$$C = B \log_2\left(1 + \text{SNR}\right) = B \log_2\left(1 + \frac{P_s}{P_n}\right) \text{ [bit/s]}$$

Donde:
\begin{itemize}
    \item $C$ = capacidad del canal (bit/s)
    \item $B$ = ancho de banda (Hz)
    \item SNR = relación señal a ruido (lineal)
\end{itemize}

En términos de OSNR óptico:
$$C = B_{opt} \log_2(1 + \text{OSNR})$$

\subsection*{2) Fórmula OSNR}
$$\text{OSNR} = P_{ch} - 10\log_{10}(N_{amp} \cdot P_{ASE,1})$$

O equivalentemente:
$$\text{OSNR}_{dB} = P_{ch,dBm} - 10\log_{10}(2 n_{sp} h \nu B_{opt}) - 10\log_{10}(N_{amp})$$

Con:
$$P_{ASE} = 2 n_{sp} h \nu (G-1) B_{opt} \approx 2 n_{sp} h \nu G B_{opt}$$
$$\text{OSNR} = \frac{P_{ch}}{N_{amp} \cdot 2 n_{sp} h \nu (G-1) B_{opt}}$$

\subsection*{3) Line rate con multiplexación de polarización}
$$R_{\text{line}} = 2 \times R_s \times \log_2(M)$$

Donde:
\begin{itemize}
    \item $R_s$ = symbol rate (Baud)
    \item $M$ = tamaño de la constelación (4 para QPSK, 16 para 16QAM, etc.)
    \item Factor 2 = multiplexación de polarización (DP)
\end{itemize}

\textbf{Ejemplo DP-64QAM a 43.7 Gbaud:}
$$R = 2 \times 43.7 \times \log_2(64) = 2 \times 43.7 \times 6 = 524.4 \text{ Gbit/s}$$

\subsection*{4) Parámetro para seleccionar transpondedor}

Se calcula la \textbf{eficiencia espectral} requerida:
$$\eta_{\text{req}} = \frac{C}{B \log_2(1 + \text{OSNR}_{EOL})}$$

Y se compara con la eficiencia del transpondedor:
$$\eta_{tr} = \frac{R_{\text{payload}}}{B_{ch}}$$

También se verifica:
$$\text{OSNR}_{EOL} \geq \text{OSNR}_{\text{min,tr}} + M$$

\section*{Problema - Tres rutas submarinas}

\subsection*{Ruta Norte (Northern)}
\textbf{Datos:}
\begin{itemize}
    \item Longitud: 3500 km
    \item $\alpha = 0.17$ dB/km
    \item Espaciado amplificadores: 65 km
    \item BW EDFA: 4000 GHz
    \item NF: 5.5 dB
    \item $P_{max} = 15$ dBm
\end{itemize}

\textbf{Cálculos:}

Número de amplificadores:
$$N_{amp} = \frac{3500}{65} \approx 54$$

Ganancia por amplificador:
$$G = \alpha \times d = 0.17 \times 65 = 11.05 \text{ dB}$$

Número de slots espectrales:
$$N_{slots} = \frac{4000}{12.5} = 320$$

\textbf{Selección transpondedor:}

Calculamos OSNR disponible con cada modelo:

\textit{Modelo 1 (DP-QPSK, 100G, 3 slots):}
$$N_{ch} = \lfloor 320/3 \rfloor = 106$$
$$P_{ch} = 15 - 10\log_{10}(106) = 15 - 20.25 = -5.25 \text{ dBm}$$
$$\text{OSNR} = -5.25 - 10\log_{10}(2 \times 10^{5.5/10} \times 1.986 \times 10^{-19} \times 194 \times 10^{12} \times 12.5 \times 10^9) - 10\log_{10}(54)$$
$$\text{OSNR} = -5.25 - (-57.48) - 17.32 = 34.91 \text{ dB}$$

Margen: $34.91 - 12 - 1 = 21.91$ dB ✓

\textit{Capacidad Shannon:}
$$C_{th} = 12.5 \times 10^9 \times \log_2(1 + 10^{34.91/10}) \times 106 = 205.2 \text{ Tbit/s}$$

\textit{Capacidad real:}
$$C_{real} = 100 \times 106 = 10.6 \text{ Tbit/s}$$

\textbf{Tabla de resultados Northern:}
\begin{center}
\begin{tabular}{|l|c|}
\hline
Spectral slot count & 320 \\
Transponder model & 1 (DP-QPSK) \\
Channel count & 106 \\
Power per channel & -5.25 dBm \\
End of line OSNR & 34.91 dB \\
Theoretical capacity & 205.2 Tbit/s \\
Actual capacity & 10.6 Tbit/s \\
\hline
\end{tabular}
\end{center}

\subsection*{Ruta Central}
\textbf{Datos:}
\begin{itemize}
    \item Longitud: 3000 km, $\alpha = 0.18$ dB/km
    \item Espaciado: 50 km, BW: 3800 GHz
    \item NF: 6 dB, $P_{max} = 15$ dBm
\end{itemize}

$$N_{amp} = 3000/50 = 60, \quad G = 0.18 \times 50 = 9 \text{ dB}$$
$$N_{slots} = 3800/12.5 = 304$$

\textbf{Modelo 1:} $N_{ch} = 101$, $P_{ch} = -5.04$ dBm, OSNR $= 33.89$ dB ✓

\textbf{Capacidad:} $C = 10.1$ Tbit/s

\subsection*{Ruta Sur (Southern)}
\textbf{Datos:}
\begin{itemize}
    \item Longitud: 4000 km, $\alpha = 0.16$ dB/km
    \item Espaciado: 75 km, BW: 4200 GHz
    \item NF: 5 dB, $P_{max} = 16$ dBm
\end{itemize}

$$N_{amp} = 4000/75 \approx 53, \quad G = 0.16 \times 75 = 12 \text{ dB}$$
$$N_{slots} = 4200/12.5 = 336$$

\textbf{Modelo 1:} $N_{ch} = 112$, $P_{ch} = -4.49$ dBm, OSNR $= 36.73$ dB ✓

\textbf{Capacidad:} $C = 11.2$ Tbit/s

\section*{Pregunta 6: Mejoras de capacidad}

\begin{enumerate}
    \item \textbf{Usar modulaciones de mayor orden:} DP-16QAM, DP-8QAM según OSNR disponible
    \item \textbf{Multiplexación espacial:} Múltiples pares de fibras (SDM)
    \item \textbf{Amplificadores de mayor potencia:} Aumentar $P_{max}$
    \item \textbf{Reducir espaciado de amplificadores:} Mejora OSNR
    \item \textbf{Técnicas de codificación:} FEC avanzado, shaping probabilístico
    \item \textbf{Usar banda L adicional:} Duplicar el espectro disponible
\end{enumerate}

\vspace{1cm}
\noindent\rule{\textwidth}{0.4pt}
\begin{center}
\textit{Fin - Examen 2022 OPTIQUE}
\end{center}

\end{document}
